\scanpage{29}
% loc: upper left, slanting upward to the right
\margin{$(\unsure{\ast\ast|\ast})$ \unsure{exprime|expriment|exprimons} que si $a\ne b$, alors $\unsure{\Pi|\pi}_a\ne\Pi_b$}
% loc: opening paragraph
De plus, si $a_1,a_2\in A$, $a_1\ne a_2$, on voit de suite que les décompositions $\Pi_{a_1},\Pi_{a_2}$ sont bien parallèles, à cela près qu'il faut exiger \deleted{\unsure{que si $s$ est un sommet d'or}}%
% loc: inserted below the cancelled continuation; continues in the left margin
\inserted{$(\ast\ast)$ si $s\in S$ est d'ordre $2$, alors un des trois ens $U_s,E_a,E_b$ est $\ne\varnothing$, où $a,b$ sont les deux arêtes inc. à $s$}

% loc: middle paragraph
Une \unsure{\deleted{façon}|façon} de visualiser la construction est via un arbre topologique $\deleted{A}\mathcal{X}=(\overset{\text{\inserted{esp. top.}}}{\underline{\Gamma}},S)$ (tel que $\Gamma$ en soit une description combinatoire). On fixe un pt sur chaque \unsure{arète|arête}, de façon à donner un graphe topologique (subdivision barycentrique) $\mathcal{X}^{\ast}=(\underline{\Gamma}^{\ast},\underbrace{S\amalg A'}_{S^{\ast}})$, où $A'\unsure{\in|\subset}\underline{\Gamma}$ est en corr \deleted{b}biunivoque donnée avec \deleted{\unsure{tout}} $A$.

% loc: bracketed NB
\deleted{pt de} [\uline{NB} $A'$ est formé de \deleted{pts} sommets d'ordre $2$, et \inserted{pour} tte arête de $\mathcal{X}^{\ast}$, une exactement des deux extrémités appartient à $A'$, l'autre à $S$ -- voici une description complète des propriétés particulières aux $\underline{\Gamma}^{\ast}$\unsure{..|...}]

% loc: lower paragraph; common letters in the article correction occur once
Alors \unsure{les|le\deleted{s}} point\deleted{s} \inserted{$a'$} de $A'$ coupe\deleted{nt} $\underline{\Gamma}^{\deleted{\ast}}$ en deux composantes connexes, dont les adhérences définissent une décomposition binaire de $\underline{\Gamma}$, d'``équateur $\{a'\}$. Deux telles décompositions \deleted{bi} de $\underline{\Gamma}$, pour $a'\ne\unsure{a''|a'}$, définissent des décompositions binaires \unsure{stric$^t$|strict$^t$} non parallèles.
% archive sheet number: 28
